\begin{savequote}
Alles in Ordnung, auch wenn Sie in Panik geraten sind und sich beschissen fühlen.
Ich kenne das. Es ist furchtbar, und es wird nicht das letzte Mal sein.
Aber vergessen Sie nicht, was Sie gesehen haben. Regel Nr. 3: Bei Herzstillstand
zuerst den \EE{eigenen} Puls fühlen.

% \begin{flushright}
\qauthor{Dickie, in: Samuel Shem: \emph{House of God.} \cite{HouseOfGodDe}}
% \end{flushright}
\end{savequote}
\chapter
[\CO{[BLS]} Basic Life Support]
{\CC{[BLS]}\\Basic Life Support}
% \AasRsN{RS VIII}
\label{chp:BLS}
\label{topic:atemstillstand-aed}
\label{topic:kreislaufstillstand-aed}
\label{topic:pediatric-life-support}

\ChapterIntro
{Die Kardiopulmonale Reanimation -- bestehend aus
Herdruckmassage, Beatmung und Elektrotherapie mittels eines
Defibrillators -- ist eine der Kernaufgaben eines Rettungs-
oder notfallmedizinischen Dienstes. Die korrekte und
routinierte Durchführung der Reanimation hat großen Einfluss
auf das Überleben des Patienten bzw. seine Lebensqualität
nach der Entlassung aus der Spitalsbehandlung.

Dieses Kapitel behandelt die Basisreanimation mittels
eines halbautomatischen Defibrillators basierend auf den
Richtlinien des European Resuscitation Council. 
}{%
\begin{description}
  \item[Maintainer:] --
  \item[Autoren:] Diverse
  \item[Reviewer:] --
  \item[Version:] {\DocVersion} ({\DocVersionInfo})
  \item[SHA1:] {\DocGitCommit}
\end{description}

{\TxTeilkapitelAASS}
}

\SectionUnderConstruction{}



% \clearpage
\begin{figure}[H]
\begin{WidePar}
\Captionof{table}{Basic Life Support mit AED.\label{tab:algorithmus-bls-aed}}
\index{Basic Life Support}\index{BLS|see{Basic Life Support}}
\includegraphics[width=\linewidth]{./extern/BLS_mit_AED.pdf}
\end{WidePar}
\end{figure}

% \clearpage
\begin{figure}[H]
\begin{WidePar}
\Captionof{table}{Pediatric Life Support.\label{tab:algorithmus-pls}}
\includegraphics[width=\linewidth]{./extern/PLS.pdf}\index{Pediatric Life
Support}\index{BLS|see{Pediatric Life Support}}
\end{WidePar}
\end{figure}

\Layout{0}{\TxBeschreibung}
{Nach der Kontrolle der Vitalfunktionen, kann der
Ersthelfer zum Ergebnis kommen, dass der Patient \EE{kein Bewusstsein} und 
\EE{keine normale Atmung} hat, d.\,h. \EE{kein Kreislauf vorhanden} ist.
In einer solchen Situation liegt die Notfalldiagnose
\T{Atem-/Kreislaufstillstand} vor.\bigskip
}{
\begin{itemize}
  \item Kein Bewusstsein
  \item Keine normale Atmung
  \item → Kein Kreislauf
\end{itemize}
}{}{}{}

\Layout{0}{Achtung}
{Entgegen der Vorgehensweise der Rettungskette wird bei einem Atem- und
Kreislaufstillstand grundsätzlich bereits \E{nach der Kontrolle der
Vitalfunktionen} der \E{Notruf abgesetzt} und erst danach mit der \E{lebensrettenden
Sofortmaßnahme} begonnen. \bigskip\bigskip
}{
\begin{itemize}
  \item Nach der Kontrolle der Vitalfunktionen Notruf absetzen.
  \item Anschließend lebensrettende Sofortmaßnahmen einleiten.
\end{itemize}
}{}{}{}

\MassnahmenNG{Erste-Hilfe-Maßnahmen}{Atem- und Kreislaufstillstand}{}{
\begin{itemize}
  \item Herz-Lungen-Wiederbelebung (Reanimation)
  \begin{itemize}
    \item Herzdruckmassage
    \item Beatmung
  \end{itemize}
\end{itemize}
}

% \Layout{0}{Lebensrettende Sofortmaßnahme
% }{
% Nach dem Absetzen des Notrufes muss sofort mit der
% \EE{Herz-Lungen-Wiederbelebung (Reanimation)} begonnen werden. Diese dient der
% Aufrechterhaltung einer von sich aus ausbleibenden Atmung und fehlenden Kreislauftätigkeit. 
% Zusätzlich zielt sie auf ein Wiederingangsetzen der eigenständigen
% Herz- sowie Lungentätigkeit ab. Die Reanimation besteht aus zwei
% Teilschritten, der \EE{Herzdruckmassage} zur Nachahmung einer
% Herztätigkeit und der \EE{Beatmung} zur Zuführung von Sauerstoff in die
% Lungen.
% }{
% \begin{itemize}
%     \item Herz-Lungen-Wiederbelebung (Reanimation)
%     \begin{itemize}
%         \item Herzdruckmassage
%         \item Beatmung
%     \end{itemize}
% \end{itemize} 
%  }

\subsubsection{Herzdruckmassage}
\index{Herzdruckmassage}

Die Herzdruckmassage bewirkt mittels Kompression des Herzens
zwischen Wirbelsäule und Brustbein ein künstliches Fließgeschehen
(Zirkulation) des Blutes.

\Layout{0}{Technik}
{
\EE{Position und Druckpunkt} Der Patient muss auf einem
harten Untergrund liegen. Ggfs. ist der Patient aus einem Bett
o.\,ä. auf den Boden zu bringen. Der Helfer kniet
dann seitlich neben dem Körper des Patienten und sucht
den richtigen \E{Druckpunkt} für die Herzdruckmassage
auf. Dieser liegt in der \EE{Mitte des Brustkorbes}.

\EE{Arm- und Körperhaltung} Die richtige
Arm- und Körperhaltung ist wichtig, um mit wenig Kraftaufwand eine
hervorragende Leistung zu erzielen. Die \E{Finger
sollten verschränkt sein}, nur der \E{Handballen liegt am Brustbein auf}.
Die \E{Arme bleiben gestreckt}, und die \E{Schultern liegen genau über
den Druckpunkt}. Nun wird mit dem Gewicht des Oberkörpers Druck auf das Brustbein ausgeübt.

\EE{Druckerzeugung} Die
Druckerzeugung hat \E{kräftig, aber nicht ruckartig} zu
erfolgen. Die \E{Eindrücktiefe} beträgt \E{\nicefrac{1}{3} des
Brustkorbes}.

\EE{Druckentlastung} Der Brustkorb muss sich wieder ganz
ausdehnen können. Es darf daher \E{kein Restdruck} bestehen. Die Druck-
und Entlastungsphasen dauern beide gleich lange.

\EE{Arbeitsgeschwindigkeit} Die Arbeitsgeschwindigkeit der
Herzdruckmassage beträgt \E{100 Massagen pro Minute}.

}{
\begin{itemize}
  \item Druckpunkt: \E{Mitte} des Brustkorbes
  \item Arm- und Körperhaltung
  \begin{itemize}
    \item Finger verschränkt,
    \item Handballen am Brustbein,
    \item Arme bleiben gestreckt,
    \item Schultern genau über den Druckpunkt.
  \end{itemize}
  \item Druckerzeugung
  \begin{itemize}
    \item kräftig aber nicht ruckartig
    \item Eindrucktiefe: \nicefrac{1}{3} des Brustkorbes
  \end{itemize}
  \item Druckentlastung
  \item 100\PerMin*
\end{itemize}
}{}{}{}

\Layout{2}{Herzdruckmassage: sooft wie möglich!}{%
Die HDM muss so rasch beginnend und so unterbrechungsfrei wie
möglich sein.

Nach Abgabe des Schocks muss die Herzdruckmassage sofort
weitergeführt werden! Auch wenn die Defibrillation
erfolgreich war, d.\,h. ein ordnungsgemäßer Herzrhythmus
wiederhergestellt wurde, kann es einige Zeit dauern, bis
wieder ein ausreichender Kreislauf einsetzt
\cite{Sandroni2008,Skhirtladze2010}.

\Fazit{Die Unterbrechung der HDM durch den AED/Defibrillator
soll insgesamt nicht mehr als \E{5\Sek*} betragen!}
}{}{}{}{}

\Layout{2}{Qualität der Herzdruckmassage}{% 
Die richtige Durchführung der Herzdruckmassage hat
wesentlichen Einfluss auf das Behandlungsergebnis bzw. das
Überleben.

Die \EE{Kompressionsrate} muss mindestens \EE{100\PerMin*}
betragen. Die \EE{Drucktiefe} soll (beim Erwachsenen)
\EE{mindestens 5\CM* und maximal 6\CM*} betragen. Damit sich
das Herz mit Blut füllen kann ist eine komplette
\EE{Entlastung} wichtig: Während der Entspannungsphase darf
der Helfer nicht auf dem Brustkorb des Patienten lehnen.
\cite[p.1283]{Erc2010Sec2}
}{}{}{}{}

\Layout{0}{Mögliche Fehlerquellen}
{Das Schlimmste, das der Ersthelfer falsch machen kann, ist
 nichts zu tun. Wenn reanimiert wird, kann es zu dem
 einen oder anderen Fehler kommen. Der Patient muss auf
 einer festen Unterlage gelagert werden, ein Bett ist
 \EE{keine feste Unterlage}! Es darf nicht vergessen werden, den Patienten auf den Boden oder eine andere harte Unterlage zu
 legen.

Grundsätzlich
kann es auch zu \EE{Verletzungen durch den Ersthelfer} kommen. Dazu zählt der
Bruch des Brustbeines durch einen zu festen Druck, Lungen- oder Herzquetschung
durch Abweichung des Druckpunktes, Rippenbrüche durch Abweichen des
Druckpunktes oder nicht senkrecht ausgeführten Druck, Organquetschung bzw.
-einreißung durch einen zu tiefen Druckpunkt.  

Ein zu ruckartiger Druck führt zu einer zu kurzen Kompression des Herzens, was
wiederum eine verkürzte Auswurfphase bedeutet. Dies alles führt zu einer
\EE{ungenügenden Zirkulation} des Blutkreislaufs. 

\Cave{Wird einmal falsch Druck ausgeübt, oder eine Rippe gebrochen, nicht 
 aufhören, sondern den
richtigen Druckpunkt erneut aufsuchen und weiter machen. 

Gebrochene Rippen und
gequetschte Organe können wieder repariert werden. Das Herz, das nicht mehr
schlägt, benötigt unsere Hilfe!}
}{
\begin{itemize}
  \item keine feste Unterlage
  \item Verletzungen durch den Ersthelfer
  \item ungenügende Zirkulation
\end{itemize}
}{}{}{}

\subsubsection{Beatmung}
\index{Beatmung!Erste-Hilfe}

\Layout{0}{Arten}
{Die bekannteste Methode zur Beatmung durch den Ersthelfer ist die
\T{Mund--zu--Mund--Beatmung}. Dabei wird bei überstrecktem Kopf mit der stirnseitigen Hand die Nase des Patienten zugedrückt. Anschließend mit dem eigenen Mund den des Patienten umschließen und
dabei die Lippen luftdicht an den Patienten drücken. Nun
wird die eigene Ausatemluft in den Patient eingeblasen. Vor der nächsten Atemspende
muss der Mund wieder abgesetzt werden, um selbst frische Luft einatmen zu
können. 

Die zweite Möglichkeit ist die \T{Mund--zu--Nase--Beatmung}. Hier wird bei überstrecktem Kopf mit der kinnseitigen Hand der Mund des Patienten
zugehalten. Anschließend mit dem eigenem Mund die Nase des Patienten umschließen und die Lippen luftdicht an den Patienten drücken. Nun wird die eigene
Ausatemluft in den Patienten eingeblasen. Vor der nächsten Atemspende
muss der Mund wieder abgesetzt werden, um selbst frische Luft einatmen zu
können. 

Bei Säuglingen und Kleinkindern wäre noch die
\T{Mund--zu--Mund/Nasen-Beatmung} möglich. In diesem
besonderen Fall wird der Mund und die Nase des Patient mit dem eigenem Mund umschlossen.
}{
\begin{itemize}
  \item Mund--zu--Mund-Beatmung
  \item Mund--zu--Nasen-Beatmung
  \item Mund--zu--Mund/Nasen-Beatmung
\end{itemize}
}{}{}{}

\Layout{0}{Wichtige Parameter für die Beatmung}
{Die \EE{Atemfrequenz} des Erwachsenen beträgt
\VonBis{12}{16}\PerMin*, d.\,h. wir atmen
\VonBis{12}{16}\PerMin* ein und aus. Das
\EE{\underline{max.} (!) Atemzugsvolumen} in ml entspricht
dem 10-fachen des Körpergewichtes, d.\,h. ein 90\Kg*
schwerer Mann kann 900\Ml* Luft ein- und wieder ausatmen.
Ein normaler Atemzug beträgt jedoch nur etwa 500\Ml* (vgl.
\See{topic:atemvolumina}).

Die ausreichende \ce{O2}-Konzentration in der Umgebungs-/Einatmungsluft ist die
Voraussetzung für eine funktionierende Atmung. Ist sie
nicht gegeben, ist auch die beste Atemarbeit nutzlos. Der normale \ce{O2}-Gehalt der Umgebungsluft beträgt
21\PC*. Der \ce{O2}-Gehalt in der \EE{Ausatemluft} beträgt
noch 16\PC*, d.\,h. der Mensch verbraucht  nur 5\PC* des
eingeatmeten Sauerstoffs. 
}{
\begin{itemize}
  \item Atemfrequenz \VonBis{12}{16}\PerMin*
  \item Max. Atemzugsvolumen [ml]: 10-fache des Körpergewichtes
  \item Einatemluft 21\PC* \ce{O2}-Gehalt
  \item Ausatemluft 16\PC* \ce{O2}-Gehalt
\end{itemize}
}{}{}{}

\Layout{0}{Bei der Beatmung ist zu beachten:}
{
Ein \EE{zu hohes Atemzugsvolumen} führt zu einem
höheren Druck im Nasen-Rachenraum, sobald die Lunge
vollständig mit Luft befüllt ist. Die Luft findet nur den
Ausweg über die Speiseröhre in den Magen. Dies kann während einer Herz-Lungen-Wiederbelebung meist unbemerkt geschehen. Bei einem \EE{zu niedrigem Atemzugsvolumen} ist die
\ce{O2}-Versorgung nicht sichergestellt, da die Luft nicht bis in die
Lunge kommt. Bei Bartträgern kann es zu einer
\EE{ungenügenden Abdichtung} kommen. Eine \EE{zu hohe
Atemfrequenz} hindert den Patienten am Ausatmen. Wenn es beim Freimachen
der Atemwege zu einer \EE{mangelnden Überstreckung} des Kopfes
kommt, kann es zu einer unzureichenden Belüftung der Lunge kommen. Bei jeder Atemspende ist der \EE{Brustkorb} des Patienten zu \EE{beobachten}. Dieser
muss sich deutlich heben und senken. Immer auch auf den \EE{Selbstschutz}
achten. Wenn möglich nur mit Beatmungstuch oder Beatmungsmaske beatmen.
}{
\begin{itemize}
  \item zu hohes Atemzugsvolumen ${\rightarrow}$ Luft in den Magen
  \item zu niedriges Atemzugsvolumen ${\rightarrow}$ \ce{O2}-Versorgung nicht
  sichergestellt
  \item ungenügende Abdichtung bei Bartträgern
  \item mangelnde Überstreckung ${\rightarrow}$ unzureichenden Belüftung der
  Lunge
  \item Brustkorb beobachten ${\rightarrow}$ deutlich heben und senken
  \item Selbstschutz ${\rightarrow}$ mit Beatmungstuch oder Beatmungsmaske
  beatmen
\end{itemize}
}{}{}{}

\subsubsection{Algorithmus Herz-Lungen-Wiederbelebung}

\Layout{0}{Vorgehensweise}{
Nach der Kontrolle des Bewusstseins und der Feststellung, dass der Patient nicht
ansprechbar ist, wird um Hilfe gerufen. Anschließend durch
Überstrecken des Kopfes die Atemwege öffnen und die Atmung kontrollieren. Wenn keine normale
Atmung vorhanden ist, Notruf absetzen. Danach mit der Herz-Lungen
Wiederbelebung beginnen. Diese setzt sich aus 30 Herzdruckmassagen und 2
Beatmungen zusammen. 
}{
\prettyref{tab:bilderserie-hlw} \ldots{} Bilderserie. 

\prettyref{tab:algorithmus-bls-aed} \ldots{} Alg. AED.
}{}{}{}

\Layout{57}{}
{}{}
{
\begin{tabularx}{\linewidth}{XXX}
\includegraphics[width=\linewidth]{./images/koch-aass/00800-gamma/Rea1b.JPG}
&
\includegraphics[width=\linewidth]{./images/koch-aass/00800-gamma/Rea2b.JPG}
&
\includegraphics[width=\linewidth]{./images/koch-aass/00800-gamma/Rea3b.JPG}
\\
% Schritt 1 
Ansprechen, berühren, Schmerzreiz
&
% Schritt 2 
Hilferuf
&
% Schritt 3 
Atemwege freimachen
\\\addlinespace
\includegraphics[width=\linewidth]{./images/koch-aass/00800-gamma/Rea4b.JPG}
&
\centering\Huge{\ttfamily 144}
&
\includegraphics[width=\linewidth]{./images/koch-aass/00800-gamma/Rea5b.JPG}
\\
% Schritt 4 
Atemkontrolle
& 
Notruf
&
% Schritt 5
30 Herzdruckmassagen
\\\addlinespace
\includegraphics[width=\linewidth]{./images/koch-aass/00800-gamma/Rea6b.JPG}
&
\includegraphics[width=\linewidth]{./images/koch-aass/00800-gamma/Rea7b.JPG}
&
\\
% Schritt 6
2 Beatmungen
&
% Schritt 7
Wenn vorhanden: Defibrillation.\ACh{Instruktor 1}{0.8.0}{Weg lassen.
}\ACh{GABS}{2010-02-19}{Wieso? Hinweis auf AED wäre doch sinnvoll.} 
&
\\\addlinespace\bottomrule
\end{tabularx}
}{Bilderserie: Herz-Lungen-Wiederbelebung
(Reanimation)\label{tab:bilderserie-hlw}\index{Bilderserie!Reanimation}\index{Reanimation!Bilderserie}}{}

% {
% {Die Durchführung der Reanimation hält sich an folgenden
% Algorithmus:}\footnote{\ {Die folgende Darstellung der
% Reanimation erlaubt nur ein Überblick über die Möglichkeiten und
% Maßnahmen der Reanimation. Eine ausführlichere Darstellung findet
% sich in den nachfolgenden Kapiteln des Lehrbuches. }}}
% 
% 
% 
% 
%  Einfügen des
% BLS-Algorithmus erstellt vom Josef 2006 für die
% ASB-921-Schulung


\Layout{0}{Ende der HLW}
{\index{Reanimation!Beendigung}
Beendet wird eine Herz-Lungenwiederbelebung durch den Ersthelfer nur
dann, wenn einer der folgenden drei Umstände eintritt. Ohne diese
zusätzlichen Umstände darf der Ersthelfer die Reanimation nicht
unterbrechen: 

\begin{description}
  \item[Eintreffen von \EE{kompetenter Hilfe.}] Der Ersthelfer darf aber beim
  Eintreffen der Rettungskräfte nicht sofort mit der Wiederbelebung  aufhören, sondern muss diese weiterführen
  bis er ausdrücklich vom Einsatzpersonal abgelöst  wird.
  \item[Erschöpfung.] Sollte der Ersthelfer all seine Kräfte aufgebraucht haben,
  wird gezwungener Maßen die Wiederbelebung ab-, bzw. unterbrochen.
  % aufgrund der
  % körperlichen Anstrengung im Zusammenhang mit der Reanimation all seine Kräfte aufbrauchen, und daher fortführen können, ist es  ihm erlaubt diese
  % einzustellen. Diese Möglichkeit wird aus den Grundsätzen der Zumutbarkeit
  % und der Versuchspflicht gewonnen und darf nur als absolut
  % letzte Option betrachtet werden.
  \item[Der Patient zeigt \EE{Lebenszeichen}.] Entweder
  erwacht der Patient und ist ansprechbar \footnote{eine
  leider äußerst seltene Form der Reanimationsbeendigung},
  oder der Patient gibt ein Husten oder Räuspern von sich
  bzw. atmet beim Beatmungsvorgang dem Ersthelfer entgegen. \EE{Erbrechen zählt \E{nicht} als Lebenszeichen.}
  % Da hier der Ersthelfer nicht mit verlässlicher Sicherheit vom Erfolg seiner  Wiederbelebung
  % ausgehen kann, muss er erneut die Atmung kontrollieren. Sollte sich das
  % Lebenszeichen durch vorhandene Atmung bestätigen, muss der Patient
  % in die stabile  Seitenlagerung gebracht und in Form der unten
  % beschriebenen anschließenden Behandlung  weiterversorgt werden.\item[]
\end{description}

}{
\begin{itemize}
  \item Eintreffen von kompetenter Hilfe
  \item Erschöpfung
  \item Patient zeigt Lebenszeichen
\end{itemize}
}{}{}{}


 

% \begin{enumerate}
%     \item  Lebenszeichen:
%     
%     {\itshape
%     Lebenszeichen können in zwei Unterkategorien auftreten:}
%     
%     \begin{enumerate}
%         \item  sichere Lebenszeichen:
%         
%         {\scriptsize
%         Der Patient erwacht und ist ansprechbar. Eine leider äußerst
%         seltene Form der  Reanimationsbeendigung.}
%         
%         \item  unsichere Lebenszeichen:
%         
%         {\scriptsize
%         Der Patient gibt ein Husten oder Räuspern von sich bzw. atmet beim
%         Beatmungsvorgang dem  Ersthelfer entgegen. Da hier der Ersthelfer nicht
%         mit verlässlicher Sicherheit vom Erfolg seiner  Wiederbelebung
%         ausgehen kann, muss er erneut die Atmung kontrollieren. Sollte sich das
%         Lebenszeichen durch vorhandene Atmung bestätigen, muss der Patient
%         in die stabile  Seitenlagerung gebracht und in Form der unten
%         beschriebenen anschließenden Behandlung  weiterversorgt werden.}
%     \end{enumerate}
%     \item  Eintreffen der Rettung:
%     
%     {\scriptsize
%     Der Ersthelfer darf bei Eintreffen der Rettungskräfte nicht sofort
%     mit der Wiederbelebung  aufhören, sondern muss diese weiterführen
%     bis er ausdrücklich vom Einsatzpersonal abgelöst  wird.}
%     
%     \item  Aufbrauchen der eigenen Kräfte:
%     
%     {\scriptsize
%     Sollte der Ersthelfer aufgrund der körperlichen Anstrengung im
%     Zusammenhang mit der  Reanimation alle seine Kräfte aufbrauchen,
%     sodass er sie nicht weiter fortführen kann, ist es  ihm erlaubt diese
%     einzustellen. Diese Möglichkeit wird aus den Grundsätzen der
%     Zumutbarkeit  und der Versuchspflicht gewonnen und darf nur als absolut
%     letzte Option verwendet werden.}
%     
% \end{enumerate}


% \subsubsection{8.1) Exkurs -- Defibrillation:}
% 
% {
% Auch für den Ersthelfer ist die Anwendung eines Defibrillators
% möglich. Die Durchführung orientiert sich hierbei an einem
% ähnlichen Ablauf wie beim Rettungssanitäter, wodurch an dieser
% Stelle nicht näher auf die Anwendung eingegangen wird. Zur näheren
% Auseinandersetzung siehe das Kapitel
% {\quotedblbase}{\dots}{\textquotedblleft}.}

\clearpage % TODO temp clearpage
\Layout{0}{Anschließende Behandlung}
{%
Die anschließenden Behandlungen orientieren sich an den
verschiedenen Beendigungsmöglichkeiten der
Reanimation, wodurch auch sie in ihrer Ausgestaltung verschieden sind. Zeigt
der Patient \EE{Lebenszeichen}, ist eine
Atemkontrolle durchzuführen und wenn eine normale Atmung
vorhanden ist, wird der Patient in die stabile Seitenlage umgelagert.
Anschließend weiterführende Maßnahmen gemäß Notfalldiagnose Bewusstlosigkeit. 

Bei Eintreffen von \EE{kompetenter Hilfe} muss der Ersthelfer
dem Rettungspersonal Auskunft über relevante Umstände und
die gesetzten Maßnahmen geben. 

Wenn der \EE{Ersthelfer erschöpft} ist, sollte er (falls überhaupt
konditionell möglich) erneut versuchen andere  Ersthelfer zu
finden. Scheitert auch dieser Versuch, kann der Ersthelfer keine weitere 
Maßnahme setzen als den Patienten zuzudecken und auf das
Eintreffen des  Rettungspersonals zu warten. Der Ersthelfer hat alles
in seiner Macht stehende getan, und es ist ihm kein Vorwurf
zu machen. 
}{
\begin{itemize}
  \item Lebenszeichen kontrollieren
  \item Atemkontrolle alle 30 Sek
  \item Stabile Seitenlage.
  \item Wärmeerhalt (zudecken)
  \item Frischluftzufuhr (Fensteröffnen, \ldots)
  \item Psychische Betreuung
  \item Kompetente Hilfe kommt ${\rightarrow}$ Auskunft über gesetzte
  Maßnahmen geben
  \item Erschöpfung: ${\rightarrow}$ versuchen weitere Ersthelfer zu
  finden
\end{itemize}

}{}{}{}




% \subsubsection{8.2) anschließende Behandlung:}
% 
% {
% Die anschließenden Behandlungen orientieren sich an den
% verschiedenen Beendigungsmöglichkeiten einer ersthelferischen
% Reanimation, wodurch auch sie in ihrer Ausgestaltung verschieden sind.}
% 
% {
% \foreignlanguage{english}{{
% a}}\foreignlanguage{english}{{d i) Lebenszeichen:}}}

% {
%  ad i) a) sichere Lebenszeichen:}
% 
% {\scriptsize
%  Der sich bei Bewusstsein befindende Patient ist in eine für ihn
% angenehme Position zu  bringen. Er ist psychisch zu betreuen und
% laufend zu beobachten. Um einen Wärmeverlust zu  vermeiden muss der
% Patient zugedeckt werden.}

% {
%  ad i) b) unsichere Lebenszeichen:}
% 
% {\scriptsize
%  Nach Umlagerung in die stabile Seitenlagerung sind die
% {\quotedblbase}ersthelferischen Allzweckwaffen{\textquotedblleft} 
% (Wärmeerhalt sichern, psychische Betreuung, laufende Kontrolle)
% durchzuführen.}

% {
%  ad ii) Eintreffen der Rettung:}
% 
% {\scriptsize
%  Der Ersthelfer muss dem Rettungspersonal Auskunft über relevante
% Umstände und gesetzte  ersthelferische Maßnahmen geben.}
% 
% {
%  ad iii) Aufbrauchen der eigenen Kräfte:}
% 
% {\scriptsize
%  Sollte der Ersthelfer sich nicht mehr körperlich in der Lage fühlen
% eine Reanimation  fortzuführen, sollte er (falls überhaupt
% konditionell möglich) erneut versuchen weitere  Ersthelfer zu finden.
% Scheiterte auch dieser Versuch, kann der Ersthelfer keine weiteren 
% Maßnahmen setzen als den Patienten zuzudecken und auf das
% Eintreffen des  Rettungspersonals zu warten. Der Ersthelfer hat alles
% in seiner Macht stehende getan und  darf sich keine Vorwürfe eines
% Versagens machen.}




% \subsection{9) Atemwegsverlegung durch Fremdkörper:}
% 
% {
% Unter dem Begriff {\quotedblbase}Bolus-Geschehen{\textquotedblleft}
% werden Fälle erfasst, bei denen das durch einen Biss (Bolus) in den
% Mundraum gelangte Objekt in die falsche Röhre, nämlich die
% Luftröhre, wandert. Bei der Behandlung muss unterschieden werden
% zwischen einer teilweisen, leichten Verlegung oder einer gänzlichen,
% schweren Verlegung.}
% 
% \subparagraph{Leichte Verlegung:}
% 
% \begin{enumerate}
%     \item  Kennzeichen:
%     \begin{enumerate}
%         \item  effektives Husten und Würgen
%     \end{enumerate}
%     \item  ersthelferische Behandlung:
%     \begin{enumerate}
%         \item  Patient beim Husten unterstützen durch kreisförmiges Reiben am
%         Rücken (andere Manipulationen sind zu unterlassen!)
%         \item Notrufabgabe
%         \item  Wärmeerhalt sichern, psychische Betreuung, laufende Kontrolle
%     \end{enumerate}
% \end{enumerate}
% 
% 
% \subparagraph{
% Schwere Verlegung:}
% Kennzeichen:
% \begin{enumerate}
%     \item  ineffektives (stummes) Husten
%     \item   Aussetzen der Atmung
%     \item   nach gewisser Zeit Bewusstlosigkeit
% 
%     \item  ersthelferische Behandlung:
% \end{enumerate}
% \begin{enumerate}
%     \item  Rückenschläge: 5x zwischen die Schulterblätter
% tangential schlagen,    danach Kontrollieren ob Objekt abgehustet
% wurde\footnote{\ {Für den Sanitäter besteht abwechselnd zu
% den Rückenschlägen, die Pflicht zur Durchführung eines
% Heimlich-Manövers. Die nähere Durchführung siehe im Kapitel
% {\quotedblbase}{\dots}{\textquotedblleft}.}}
% 
%     \item  wurde das Objekt nicht abgehustet und ist der Patient weiterhin bei
%    Bewusstsein: erneute 5 Rückenschläge mit anschließender
% Kontrolle
% 
%     \item  wurde das Objekt nicht abgehustet und ist der Patient bewusstlos
% (und   zwingend ohne Atmung und Kreislauf):    
% 
% {
%   Notrufabgabe + Wiederbelebung (auch hier ist die Beatmung zwingend   
% durchzuführen; es besteht die Chance durch die Beatmung das Objekt in
% den   rechten Bronchus zu bewegen)}
% 
%     \item  wurde das Objekt abgehustet und ist der Patient bei Bewusstsein: 
% 
% {
%   Lagerung mit erhöhten Oberkörper, Notrufabgabe, Wärmeerhalt
% sichern,   psychische Betreuung, laufende Kontrolle}
% 
%     \item wurde das Objekt abgehustet und ist der Patient bewusstlos:
% 
% {
%   bei aufrechter Atmung/Kreislauf: stabile Seitenlagerung, Notrufabgabe,
%    Wärmeerhalt sichern, psychische Betreuung, laufende Kontrolle}
% 
% {
%   bei fehlen von Atmung/Kreislauf: Notrufabgabe + Wiederbelebung}
% \end{enumerate}
% {
% Selbst bei relativ schneller Abhustung des Fremdkörpers, ist es ratsam
% eine ärztliche Kontrolle aufzusuchen, da es im Zuge eines
% Bolus-Geschehn leicht zu {\quotedblbase}hypoxischen
% Schäden{\textquotedblleft} (Schäden durch Sauerstoffmangel) kommen
% kann.}


\section{Anmerkungen}



\Layout{2}{Infektionen}{% 
Die Ansteckungsgefahr mit Krankheiten im Rahmen einer
Reanimation sind sehr gering. Es wurden noch keine
Infektionen mit schwerwiegenden Erkrankungen wie HIV oder
Hepatitis C während einer Basisreanimation beschrieben.
\cite[p.1284]{Erc2010Sec2} }{}{}{}{}

\Layout{2}{Probleme bei der HDM während des Defi-Ladevorganges}{% 
\Ca{Der folgende Kommentar betrifft nur Helfer, welche nach
den ERC~2010-Leitlinien arbeiten!} 
In den ERC-Leitlinien von
2010 wird erstmals empfohlen, die Herzdruckmassage fortzusetzen,
\EE{während der Defibrillator sich auflädt}, d.\,h. nach der
Analyse wird weiter massiert, bis der Defi bereit zur
Schockabgabe ist. Dies soll die HDM-Pausen verringern. Bei
vielen AED-Geräten kann es jedoch zu
Problemen kommen, sie interpretieren die
reanimationsbedingten Störungen der EKG-Ableitung als
Kontaktverlust bzw. Bewegungsartefakt und brechen den
Ladevorgang ab.\footnote{Davon betroffen sind derzeit z.\,B. die
Geräte der Schiller Fred{\texttrademark}-Serie, der
Defigard{\texttrademark} 1002 und der
LifePak{\texttrademark} 500. Bezüglich anderer Geräte liegen
der Redaktion derzeit noch keine Informationen vor. Es muss
jedoch davon ausgegangen werden, dass auch andere Modelle
ein derartiges Verhalten zeigen können. \E{(Angaben ohne
Gewähr. Stand: 2011-01-19)}} Dadurch wird die Schockabgabe
verhindert. Daher kann diese Methode nicht generell
empfohlen werden, sondern darf nur durchgeführt werden, wenn
der Hersteller des Gerätes dies freigegeben hat.
\Commentary[Herzmassage während AED-Ladezyklus]{Lt.
vorläufiger Auskunft der Distributoren ist dies momentan
zumindest mit den AED-Geräten der Fa. Schiller
(Fred{\texttrademark}, Defigard{\texttrademark} 1002) nicht
möglich, da auch während der Aufladephase das EKG über die
Defi-Elektroden überwacht wird. Die elektrische Aktivität
durch Thoraxkompressionen wird dann als Bewegungsartefakt
interpretiert und es wird eine interne Sicherheitsentladung
eingeleitet. Für die Modelle der Fred-Reihe läuft momentan
eine Evaluierung bezüglich eines Softwareupdates, es wird
daher aller Voraussicht nach in Zukunft eine Lösung geben.
 
Für die Geräte der Corpuls{\texttrademark}-Reihe liegen
uns derzeit noch keine Informationen vor.  
Der LifePak{\texttrademark} 500 von Medtronic verhält sich
wie der Defigard{\texttrademark} 1002. Bezüglich der Typen
LifePak{\texttrademark} 1000, 12 und 15 liegen ebanfalls
noch keine Informationen vor. 
Es muss
 davon ausgegangen werden, dass auch andere Geräte
ein derartiges Verhalten zeigen. \E{(Angaben ohne
Gewähr. Stand: 2011-01-19)}}

\Cave{Manche Defibrillatoren brechen den Ladevorgang ab,
wenn währenddessen eine Herzdruckmassage durchgeführt wird.}

\Fazit{Während des Ladevorganges des Defibrillators darf nur
bei Verwendung \E{dafür} vorgesehener Geräte eine Herzdruckmassage
durchgeführt werden.} 
}{}{}{}{}
